\section{总结与延伸}

\subsection{一条主线：知识存在，不等于知识可用}

本讲最重要的概念不是某一个 reversal benchmark，而是把“learning”拆成四层：经验提供 explicit content；内容蕴含 latent information；训练过程形成内部表示；测试时接口决定其中多少成为 usable information。参数与上下文的差异，正发生在后两层。

\begin{importantbox}{最终知识框架}
\begin{enumerate}
  \item Parametric learning 与 ICL 可以吸收同一事实，却形成不同的访问与泛化行为。
  \item Reversal、syllogism 与 codebook 是测试 latent relation 的受控探针，不是所有智能的完整代理。
  \item Co-occurrence 能提高真实性能，也能掩盖 systematic inference 的缺失。
  \item In-context augmentation 把 latent implication 离线变成 explicit supervision。
  \item Episodic retrieval 在目标出现时把原经历恢复到 context，但真实 recall 仍是难题。
  \item RL regeneration 学习在 test time 重建信息，对部分结构迁移，对 true reversal 仍受搜索限制。
  \item 三条路线应按训练成本、推理成本、覆盖、可追溯性和失败恢复组合使用。
  \item 自然智能类比支持互补记忆系统这一计算需求，不证明神经实现同一。
\end{enumerate}
\end{importantbox}

\subsection{从课堂结果到系统架构}

一个实际系统可以采用分层策略：对高频、稳定、可验证的 latent relation 做离线 augmentation；对长尾事实和用户经历使用可追溯 retrieval；对无法直接检索、但具有通用结构的任务使用 bounded test-time reasoning；当检索置信度低或推理搜索过长时，触发澄清、工具调用或人工升级。

\begin{lstlisting}[language=Python,caption={组合三种 bridge 的部署决策策略}]
def answer(query):
    direct = parametric_model.answer(query)
    if verifier.confident(direct):
        return direct

    memories = retriever.search(query)
    if retriever.has_high_recall_evidence(memories):
        return grounded_model.answer(query, context=memories)

    reasoning = reasoning_policy.solve(query, token_budget=budget(query))
    if verifier.confident(reasoning.answer):
        return reasoning.answer

    return request_clarification_or_human_review(query)
\end{lstlisting}

\begin{knowledgebox}{建议的评测矩阵}
不要只测 explicit QA。至少交叉测试 forward、reversal、multi-hop、alternative goal、cross-lingual、long-context distractor、retrieval miss、conflicting memory 与 compute-limited reasoning，并分别报告 accuracy、latency、token、recall、provenance completeness 与 calibration。
\end{knowledgebox}

\subsection{仍然开放的研究问题}

第一，怎样训练参数表示，使其天然支持多方向索引，而不靠枚举？第二，能否让模型学会决定“该写入参数、该保留 episode、还是只需短期 context”？第三，如何训练通用 retriever，使 recall 不依赖已知任务模板？第四，RL reasoning 能否学习检索式内部操作，同时限制 token 和错误自信？第五，怎样把 augmentation 的 prompt leakage、hallucination 与 drift 变成可审计指标？

\begin{warningbox}{三个不能从本讲推出的结论}
不能推出 ICL 总是优于 fine-tuning；不能推出所有参数知识都无法做关系推理；不能推出 hippocampus 等同于 RAG。课程证明的是特定受控条件下的信息可用性差异，并提出三类可检验、可组合的桥接方法。
\end{warningbox}

\subsection{拓展阅读}

建议按证据链阅读：先读 Berglund 等人的 Reversal Curse，理解 directional access failure；再读 Lampinen 等人的 controlled ICL/fine-tuning study，核对整数据集 context、syllogism 与 from-scratch controls；随后读 Latent Learning，重点看 episodic oracle 与 complementary memory；最后读 Improving Latent Generalization Using Test-time Compute，检查 RL transfer 与 reversal limitation。关于 ICL 与优化的关系，可对照 Akyürek 等人与 Dai 等人的受限理论结果。

\begin{knowledgebox}{Lecture snapshot}
本讲发生于 2026-05-07。讲义使用当时已公开的 arXiv 版本：controlled study v3、Latent Learning v3、test-time compute v1、Reversal Curse v4。后续 revision 可以作为新证据，但不应倒推成课堂当日已经展示的结果。
\end{knowledgebox}

最终，这堂课把“模型记住了多少”改写为一个更精确的系统问题：\textbf{当目标变化时，模型能否把过去经验重新构造成当前可计算的信息？}参数、context、episodic memory、augmentation 与 RL 并不是五个孤立模块，而是对同一信息生命周期的不同分工。

\end{document}
